% GENERATED FILE. Edit canonical lessons, then run npm run generate:books.
% canonical-source-hash: fnv1a64:1d4518256bc2c7ff
% canonical-lessons: KA-C447-rastetade, KA-C447-teravugolisu, KA-C447-addipadisu, KA-C447-sthagitagollu, KA-C447-sancarisu, KA-R447-review-441-447

\chapter{A Blocked Road}
\label{ch:ka-a-blocked-road}
\hlchaptermodality{447}

\begin{chapteropening}
\textbf{By the end of this chapter:} \emph{I can say a roadblock, to clear away, to obstruct, to stall and to ply.}

Complete the last lesson of chapter 447: I can say a roadblock, to clear away, to obstruct, to stall and to ply.
\end{chapteropening}

\section[\texorpdfstring{\kn{ರಸ್ತೆತಡೆ} (rastetaḍe)}{ರಸ್ತೆತಡೆ (rastetaḍe)}]{\kn{ರಸ್ತೆತಡೆ} (rastetaḍe) — a roadblock}

\label{lesson:KA-C447-rastetade}



\begin{quote}
Before the new one: say the Kannada for to lose one's way, then the Kannada for a footstep.
\end{quote}

\subsection*{You'll want to know: \kn{ರಸ್ತೆತಡೆ}}
\textbf{\kn{ರಸ್ತೆತಡೆ}} — \emph{rastetaḍe} — \textquotedblleft{}a roadblock\textquotedblright{}.

\begin{quote}
\kn{ಮುಷ್ಕರದ} \kn{ಕಾರಣ} \kn{ಹೆದ್ದಾರಿಯಲ್ಲಿ} \kn{ರಸ್ತೆತಡೆ} \kn{ಇತ್ತು}.
\end{quote}

\begin{quote}
\emph{muṣkarada kāraṇa heddāriyalli rastetaḍe ittu.}
\end{quote}

\begin{quote}
Because of the strike there was a roadblock on the highway.
\end{quote}

\kn{ರಸ್ತೆ}, road, plus \kn{ತಡೆ}, a stop: a road held shut.

\begin{cousinweb}
\textbf{\kn{ರಸ್ತೆತಡೆ}} (\emph{rastetaḍe}) is a Kannada compound of \kn{ರಸ್ತೆ}, \textquotedblleft{}a road\textquotedblright{}, usually traced to Persian \emph{rāsta}, and native \kn{ತಡೆ}, \textquotedblleft{}to stop, to block\textquotedblright{} (DEDR 3031).
\end{cousinweb}

\begin{grammarlens}[title={What you've built}]
One more word for getting around.
\end{grammarlens}

\subsection*{Your turn}
\begin{itemize}
\raggedright
  \item \emph{Say it:} \emph{rastetaḍe}
  \item \emph{Say it:} \emph{rastetaḍe}, once more
  \item \emph{Say it:} \emph{muṣkarada kāraṇa heddāriyalli rastetaḍe ittu.} — Because of the strike there was a roadblock on the highway.
  \item \emph{Recall:} say the Kannada for to lose one's way, then the Kannada for a footstep, then say \emph{rastetaḍe} again
\end{itemize}

\begin{tcolorbox}[breakable,colback=teal!4,colframe=teal!35!black,title={Before you move on}]
What does \kn{ರಸ್ತೆತಡೆ} mean? (\textquotedblleft{}A roadblock\textquotedblright{}.) Say it once more.
\end{tcolorbox}
\section[\texorpdfstring{\kn{ತೆರವುಗೊಳಿಸು} (teravugoḷisu)}{ತೆರವುಗೊಳಿಸು (teravugoḷisu)}]{\kn{ತೆರವುಗೊಳಿಸು} (teravugoḷisu) — to clear away}

\label{lesson:KA-C447-teravugolisu}



\begin{quote}
Before the new one: say the Kannada for a footstep, then the Kannada for a roadblock.
\end{quote}

\subsection*{You'll want to know: \kn{ತೆರವುಗೊಳಿಸು}}
\textbf{\kn{ತೆರವುಗೊಳಿಸು}} — \emph{teravugoḷisu} — \textquotedblleft{}to clear away\textquotedblright{}.

\begin{quote}
\kn{ಕೆಲಸಗಾರರು} \kn{ರಸ್ತೆಯಿಂದ} \kn{ಮರವನ್ನು} \kn{ತೆರವುಗೊಳಿಸಿದರು}.
\end{quote}

\begin{quote}
\emph{kelasagāraru rasteyinda maravannu teravugoḷisidaru.}
\end{quote}

\begin{quote}
The workers cleared the tree from the road.
\end{quote}

Past (I …): \kn{ತೆರವುಗೊಳಿಸಿದೆ}, \emph{teravugoḷiside}.

\kn{ತೆರವು} is an opening; -\kn{ಗೊಳಿಸು} makes it a verb: to make open.

\begin{cousinweb}
\textbf{\kn{ತೆರವುಗೊಳಿಸು}} (\emph{teravugoḷisu}) is native: \emph{teṟa}, \emph{teṟavu}, \textquotedblleft{}an opening, a clearing\textquotedblright{} (DEDR 3259), with the Kannada verb-maker -\emph{goḷisu}.
\end{cousinweb}

\begin{grammarlens}[title={What you've built}]
One more word for getting around.
\end{grammarlens}

\subsection*{Your turn}
\begin{itemize}
\raggedright
  \item \emph{Say it:} \emph{teravugoḷisu}
  \item \emph{Say it:} \emph{teravugoḷisu}, once more
  \item \emph{Say it:} \emph{kelasagāraru rasteyinda maravannu teravugoḷisidaru.} — The workers cleared the tree from the road.
  \item \emph{Recall:} say the Kannada for a footstep, then the Kannada for a roadblock, then say \emph{teravugoḷisu} again
\end{itemize}

\begin{tcolorbox}[breakable,colback=teal!4,colframe=teal!35!black,title={Before you move on}]
What does \kn{ತೆರವುಗೊಳಿಸು} mean? (\textquotedblleft{}To clear away\textquotedblright{}.) Say it once more.
\end{tcolorbox}
\section[\texorpdfstring{\kn{ಅಡ್ಡಿಪಡಿಸು} (aḍḍipaḍisu)}{ಅಡ್ಡಿಪಡಿಸು (aḍḍipaḍisu)}]{\kn{ಅಡ್ಡಿಪಡಿಸು} (aḍḍipaḍisu) — to obstruct; to disrupt}

\label{lesson:KA-C447-addipadisu}



\begin{quote}
Before the new one: say the Kannada for a roadblock, then the Kannada for to clear away.
\end{quote}

\subsection*{You'll want to know: \kn{ಅಡ್ಡಿಪಡಿಸು}}
\textbf{\kn{ಅಡ್ಡಿಪಡಿಸು}} — \emph{aḍḍipaḍisu} — \textquotedblleft{}to obstruct; to disrupt\textquotedblright{}.

\begin{quote}
\kn{ಬಿದ್ದ} \kn{ಮರ} \kn{ಸಂಚಾರವನ್ನು} \kn{ಅಡ್ಡಿಪಡಿಸಿತು}.
\end{quote}

\begin{quote}
\emph{bidda mara sañcāravannu aḍḍipaḍisitu.}
\end{quote}

\begin{quote}
The fallen tree held up the traffic.
\end{quote}

Past (I …): \kn{ಅಡ್ಡಿಪಡಿಸಿದೆ}, \emph{aḍḍipaḍiside}.

\kn{ಅಡ್ಡಿ} is a hindrance: \kn{ಅಡ್ಡಿ} \kn{ಇಲ್ಲ}, there is no objection.

\begin{cousinweb}
\textbf{\kn{ಅಡ್ಡಿಪಡಿಸು}} (\emph{aḍḍipaḍisu}) is native: DEDR 83 lists Kannada \emph{aḍḍi}, \textquotedblleft{}an obstacle, opposition, delay\textquotedblright{}, here with the verb-maker -\emph{paḍisu}.
\end{cousinweb}

\begin{grammarlens}[title={What you've built}]
One more word for getting around.
\end{grammarlens}

\subsection*{Your turn}
\begin{itemize}
\raggedright
  \item \emph{Say it:} \emph{aḍḍipaḍisu}
  \item \emph{Say it:} \emph{aḍḍipaḍisu}, once more
  \item \emph{Say it:} \emph{bidda mara sañcāravannu aḍḍipaḍisitu.} — The fallen tree held up the traffic.
  \item \emph{Recall:} say the Kannada for a roadblock, then the Kannada for to clear away, then say \emph{aḍḍipaḍisu} again
\end{itemize}

\begin{tcolorbox}[breakable,colback=teal!4,colframe=teal!35!black,title={Before you move on}]
What does \kn{ಅಡ್ಡಿಪಡಿಸು} mean? (\textquotedblleft{}To obstruct; to disrupt\textquotedblright{}.) Say it once more.
\end{tcolorbox}
\section[\texorpdfstring{\kn{ಸ್ಥಗಿತಗೊಳ್ಳು} (sthagitagoḷḷu)}{ಸ್ಥಗಿತಗೊಳ್ಳು (sthagitagoḷḷu)}]{\kn{ಸ್ಥಗಿತಗೊಳ್ಳು} (sthagitagoḷḷu) — to stall; to come to a standstill}

\label{lesson:KA-C447-sthagitagollu}



\begin{quote}
Before the new one: say the Kannada for to clear away, then the Kannada for to obstruct.
\end{quote}

\subsection*{You'll want to know: \kn{ಸ್ಥಗಿತಗೊಳ್ಳು}}
\textbf{\kn{ಸ್ಥಗಿತಗೊಳ್ಳು}} — \emph{sthagitagoḷḷu} — \textquotedblleft{}to stall; to come to a standstill\textquotedblright{}.

\begin{quote}
\kn{ಮುಷ್ಕರದಿಂದ} \kn{ನಗರದ} \kn{ಎಲ್ಲಾ} \kn{ಬಸ್ಸುಗಳು} \kn{ಸ್ಥಗಿತಗೊಂಡವು}.
\end{quote}

\begin{quote}
\emph{muṣkaradinda nagarada ellā bassugaḷu sthagitagoṇḍavu.}
\end{quote}

\begin{quote}
Because of the strike all the city buses came to a standstill.
\end{quote}

Past (it …): \kn{ಸ್ಥಗಿತಗೊಂಡಿತು}, \emph{sthagitagoṇḍitu}.

Said of services, work and traffic that stop.

\begin{cousinweb}
\textbf{\kn{ಸ್ಥಗಿತಗೊಳ್ಳು}} (\emph{sthagitagoḷḷu}) is built from Sanskrit \emph{sthagita}, \textquotedblleft{}covered, closed, shut\textquotedblright{} (Monier-Williams), and the Kannada verb-maker -\emph{goḷḷu}: to become shut.
\end{cousinweb}

\begin{grammarlens}[title={What you've built}]
One more word for getting around.
\end{grammarlens}

\subsection*{Your turn}
\begin{itemize}
\raggedright
  \item \emph{Say it:} \emph{sthagitagoḷḷu}
  \item \emph{Say it:} \emph{sthagitagoḷḷu}, once more
  \item \emph{Say it:} \emph{muṣkaradinda nagarada ellā bassugaḷu sthagitagoṇḍavu.} — Because of the strike all the city buses came to a standstill.
  \item \emph{Recall:} say the Kannada for to clear away, then the Kannada for to obstruct, then say \emph{sthagitagoḷḷu} again
\end{itemize}

\begin{tcolorbox}[breakable,colback=teal!4,colframe=teal!35!black,title={Before you move on}]
What does \kn{ಸ್ಥಗಿತಗೊಳ್ಳು} mean? (\textquotedblleft{}To stall; to come to a standstill\textquotedblright{}.) Say it once more.
\end{tcolorbox}
\section[\texorpdfstring{\kn{ಸಂಚರಿಸು} (sañcarisu)}{ಸಂಚರಿಸು (sañcarisu)}]{\kn{ಸಂಚರಿಸು} (sañcarisu) — to ply; to travel about}

\label{lesson:KA-C447-sancarisu}



\begin{quote}
Before the new one: say the Kannada for to obstruct, then the Kannada for to stall.
\end{quote}

\subsection*{You'll want to know: \kn{ಸಂಚರಿಸು}}
\textbf{\kn{ಸಂಚರಿಸು}} — \emph{sañcarisu} — \textquotedblleft{}to ply; to travel about\textquotedblright{}.

\begin{quote}
\kn{ಈ} \kn{ಮಾರ್ಗದಲ್ಲಿ} \kn{ಪ್ರತಿ} \kn{ಹತ್ತು} \kn{ನಿಮಿಷಕ್ಕೆ} \kn{ಬಸ್} \kn{ಸಂಚರಿಸುತ್ತದೆ}.
\end{quote}

\begin{quote}
\emph{ī mārgadalli prati hattu nimiṣakke bas sañcarisuttade.}
\end{quote}

\begin{quote}
A bus runs on this route every ten minutes.
\end{quote}

Past (I …): \kn{ಸಂಚರಿಸಿದೆ}, \emph{sañcariside}.

Said of buses and trains that run a route, and of people who travel about.

\begin{cousinweb}
\textbf{\kn{ಸಂಚರಿಸು}} (\emph{saṃcarisu}) is built from Sanskrit \emph{sam-car}, \textquotedblleft{}to move about\textquotedblright{} (the noun \emph{saṃcāra}, CDIAL 12869), with the Kannada verb ending -\emph{isu}.
\end{cousinweb}

\begin{grammarlens}[title={What you've built}]
One more word for getting around.
\end{grammarlens}

\subsection*{Your turn}
\begin{itemize}
\raggedright
  \item \emph{Say it:} \emph{sañcarisu}
  \item \emph{Say it:} \emph{sañcarisu}, once more
  \item \emph{Say it:} \emph{ī mārgadalli prati hattu nimiṣakke bas sañcarisuttade.} — A bus runs on this route every ten minutes.
  \item \emph{Recall:} say the Kannada for to obstruct, then the Kannada for to stall, then say \emph{sañcarisu} again
\end{itemize}

\begin{tcolorbox}[breakable,colback=teal!4,colframe=teal!35!black,title={Before you move on}]
What does \kn{ಸಂಚರಿಸು} mean? (\textquotedblleft{}To ply; to travel about\textquotedblright{}.) Say it once more.
\end{tcolorbox}
\section[(dialogue)]{Looking back: from Paying Charges to A Blocked Road}

\label{lesson:KA-R447-review-441-447}



\begin{quote}
Say the Kannada for a turn in line and to waive.
\end{quote}

\subsection*{Your turn}
Say these aloud:

\begin{itemize}
\raggedright
  \item a turn in line, to waive, domestic, cash, a card
  \item a highway, a toll, a flyover, a mountain pass, a route
  \item traffic, a traffic signal, a pedestrian, congestion, to control
  \item to comply with, to violate, to obey, to ban, one-way
  \item to park a vehicle, to turn round, to wobble, to move, to approach
  \item a footpath, a shortcut, a detour, to lose one's way, a footstep
  \item a roadblock, to clear away, to obstruct, to stall, to ply
\end{itemize}

\begin{tcolorbox}[breakable,colback=teal!4,colframe=teal!35!black,title={Before you move on}]
A turn in line? (\textbf{\kn{ಸರತಿ}}, \emph{sarati}.) To waive? (\textbf{\kn{ಮನ್ನಾ} \kn{ಮಾಡು}}, \emph{mannā māḍu}.) Domestic? (\textbf{\kn{ದೇಶೀಯ}}, \emph{dēśīya}.) Cash? (\textbf{\kn{ನಗದು}}, \emph{nagadu}.) A card? (\textbf{\kn{ಕಾರ್ಡ್}}, \emph{kārḍ}.) A highway? (\textbf{\kn{ಹೆದ್ದಾರಿ}}, \emph{heddāri}.) A toll? (\textbf{\kn{ಟೋಲ್}}, \emph{ṭōl}.) A flyover? (\textbf{\kn{ಮೇಲ್ಸೇತುವೆ}}, \emph{mēlsētuve}.) A mountain pass? (\textbf{\kn{ಘಾಟಿ}}, \emph{ghāṭi}.) A route? (\textbf{\kn{ಮಾರ್ಗ}}, \emph{mārga}.) Traffic? (\textbf{\kn{ಸಂಚಾರ}}, \emph{sañcāra}.) A traffic signal? (\textbf{\kn{ಸಿಗ್ನಲ್}}, \emph{signal}.) A pedestrian? (\textbf{\kn{ಪಾದಚಾರಿ}}, \emph{pādacāri}.) Congestion? (\textbf{\kn{ದಟ್ಟಣೆ}}, \emph{daṭṭaṇe}.) To control? (\textbf{\kn{ನಿಯಂತ್ರಿಸು}}, \emph{niyantrisu}.) To comply with? (\textbf{\kn{ಅನುಸರಿಸು}}, \emph{anusarisu}.) To violate? (\textbf{\kn{ಉಲ್ಲಂಘಿಸು}}, \emph{ullaṅghisu}.) To obey? (\textbf{\kn{ಪಾಲಿಸು}}, \emph{pālisu}.) To ban? (\textbf{\kn{ನಿಷೇಧಿಸು}}, \emph{niṣēdhisu}.) One-way? (\textbf{\kn{ಏಕಮುಖ}}, \emph{ēkamukha}.) To park a vehicle? (\textbf{\kn{ಪಾರ್ಕ್} \kn{ಮಾಡು}}, \emph{pārk māḍu}.) To turn round? (\textbf{\kn{ತಿರುಗು}}, \emph{tirugu}.) To wobble? (\textbf{\kn{ಅಲುಗಾಡು}}, \emph{alugāḍu}.) To move? (\textbf{\kn{ಚಲಿಸು}}, \emph{calisu}.) To approach? (\textbf{\kn{ಸಮೀಪಿಸು}}, \emph{samīpisu}.) A footpath? (\textbf{\kn{ಕಾಲುದಾರಿ}}, \emph{kāludāri}.) A shortcut? (\textbf{\kn{ಒಳದಾರಿ}}, \emph{oḷadāri}.) A detour? (\textbf{\kn{ಸುತ್ತುದಾರಿ}}, \emph{suttudāri}.) To lose one's way? (\textbf{\kn{ದಾರಿ} \kn{ತಪ್ಪು}}, \emph{dāri tappu}.) A footstep? (\textbf{\kn{ಹೆಜ್ಜೆ}}, \emph{hejje}.) A roadblock? (\textbf{\kn{ರಸ್ತೆತಡೆ}}, \emph{rastetaḍe}.) To clear away? (\textbf{\kn{ತೆರವುಗೊಳಿಸು}}, \emph{teravugoḷisu}.) To obstruct? (\textbf{\kn{ಅಡ್ಡಿಪಡಿಸು}}, \emph{aḍḍipaḍisu}.) To stall? (\textbf{\kn{ಸ್ಥಗಿತಗೊಳ್ಳು}}, \emph{sthagitagoḷḷu}.) To ply? (\textbf{\kn{ಸಂಚರಿಸು}}, \emph{sañcarisu}.)
\end{tcolorbox}
